%% parked.tex -- Parked, Not Idle (extended)
\documentclass[sigconf,nonacm]{acmart}
\microtypesetup{expansion=false}
\usepackage{xspace,booktabs,amsmath,amssymb,graphicx}
\newcommand{\sys}{\textsc{ParkAware}\xspace}
\newcommand{\code}[1]{\texttt{\small #1}}

\begin{document}

\title{Parked, Not Idle: Why Autoscaling Breaks on\\Agentic LLM Workloads}
\author{Anonymous Author(s)}
\affiliation{\institution{Submission \#xxx}\country{}}

\begin{abstract}
Every LLM autoscaler in production --- Kubernetes HPA on GPU utilisation,
KEDA on queue depth, llm-d's KV-cache-driven Workload Variant Autoscaler
--- inherits an assumption from stateless web serving: a request that is
not computing is not occupying the system. It has either finished, and
released its resources, or it is queued, and therefore visible.

Agentic workloads violate this. An agent program is a sequence of model
invocations separated by tool executions, and during a tool call the
program consumes no compute while holding its entire accumulated KV cache
resident. It is \emph{parked}, not idle: invisible to every compute-side
signal, yet fully occupying memory.

We quantify the consequences with a trace-driven study built on the public
Azure LLM inference traces. Azure's coding workload, the closest
production analogue of an agent, differs from chat in kind rather than
degree: index of dispersion 71.3 against 4.0, peak-to-mean 12.8 against
1.8, context-to-generation ratio 73:1 against 5.5:1. Simulating an elastic
cluster, we find that at eight turns per program \textbf{48.6\% of
resident programs are parked} at any instant.

The effect on autoscaling is severe. Across the agentic region of our
sweep the deployed policies attain 0.25--0.37 of the SLO and complete
0.17--0.31 of offered programs, against 0.66--0.99 SLO on the same trace
replayed as chat. We evaluate six policies spanning all three branches of
the autoscaling taxonomy, including a proactive forecaster and a tabular
Q-learner. The Q-learner is the worst of all, completing \emph{zero}
programs, because a 180-second cold start makes the credit assignment
between an immediate cost saving and a delayed SLO collapse essentially
unlearnable.

\sys changes one thing --- occupancy counts parked programs --- and
attains 0.655 SLO and 0.527 goodput, $1.46\times$ and $1.28\times$ the
best baseline. Its margin is a direct function of how expensive state is
to rebuild: at a 30\,s cold start it leads by $1.04\times$, at a realistic
180\,s for a 70B model by $2.48\times$ over HPA.
\end{abstract}

\keywords{LLM serving, agents, autoscaling, elasticity, KV cache}
\maketitle

% =====================================================================
\section{Introduction}

An autoscaler answers one question: how many replicas should be running
right now? For two decades the answer came from a utilisation signal, and
underneath every such signal lies the same assumption. Work that is not
being executed is not present. It has finished and released its
resources, or it is waiting and shows up in a queue. Either way, a
utilisation reading of zero means an empty machine.

LLM agents falsify this. An agent program is not a request; it is a loop.
The model emits an action, an external tool executes it, the result is
appended to the transcript, and the model is invoked again on the extended
context~\cite{react}. Between invocations the program computes nothing.
It also releases nothing: its KV cache, now holding the whole transcript,
stays resident, because evicting it means re-prefilling tens of thousands
of tokens when the tool returns. The program is \emph{parked}.

A parked program is invisible to every compute-side signal an autoscaler
can read. \code{DCGM\_FI\_DEV\_GPU\_UTIL} reports SM occupancy and sees
nothing, because no kernel runs on its behalf.
\code{vllm:num\_requests\_waiting} counts requests queued for a replica
and sees nothing, because the program is not waiting for a replica --- it
is waiting for a shell command to return. To an autoscaler reading those
signals, a cluster full of parked agents looks like a cluster that should
be shrunk.

Shrinking it is worse than merely wrong. Removing a replica destroys the
KV cache of every program parked on it, and each of those programs must
re-prefill its full transcript when its tool returns. That recompute
lands on a cluster that was just made smaller, driving latency up, which
the autoscaler eventually notices --- and then pays a cold start of
minutes to undo a decision it made in seconds. The asymmetry between how
fast you can destroy state and how slowly you can rebuild it is the whole
problem.

\paragraph{This paper.} We ask what happens to autoscaling when the
workload becomes agentic, and answer with a trace-driven simulation study.

\begin{itemize}
\item \textbf{Characterisation.} Using the public Azure LLM inference
  traces~\cite{splitwise}, we show the coding workload differs from chat
  in kind: $18\times$ the arrival burstiness and $13\times$ the
  context-to-generation ratio. The workload that most needs an autoscaler
  is the one on which autoscalers fail (\S\ref{sec:motivation}).

\item \textbf{Quantifying parking.} At eight turns per program, 48.6\% of
  resident programs are inside a tool call at any instant, holding memory
  and consuming no compute (\S\ref{sec:eval-parked}).

\item \textbf{A broad evaluation.} We implement six autoscalers spanning
  all three branches of the standard taxonomy --- reactive (HPA, KEDA,
  KV-util), proactive (Holt forecasting), and learning-based (tabular
  Q-learning, with and without the parked signal) --- and show all of them
  degrade sharply, in different ways (\S\ref{sec:eval-main}).

\item \textbf{A negative result on RL.} The Q-learner is the worst policy
  we tested, completing zero programs. The failure is structural, not a
  tuning artefact: cost is immediate and legible in the reward, SLO damage
  is delayed by a full cold start, and the agent reliably learns to starve
  the cluster (\S\ref{sec:eval-rl}).

\item \textbf{A mechanism, and evidence that it is the mechanism.} \sys
  counts parked programs as occupancy and recovers most of the loss. Its
  advantage is a monotone function of cold-start cost, vanishing when
  rebuilding state is cheap --- which is what we would expect if, and only
  if, the mechanism we claim is the one operating
  (\S\ref{sec:eval-cold}).
\end{itemize}

\paragraph{Scope.} This is simulation driven by real production arrival
traces. We did not run vLLM on GPUs, and \S\ref{sec:threats} states
precisely what that does and does not license us to claim.

% =====================================================================
\section{Background and Motivation}
\label{sec:motivation}

\subsection{Agentic serving}

A chat request is a function call: prompt in, completion out, state
discarded. An agent program is a stateful loop. In the ReAct
pattern~\cite{react} that underlies most deployed agents, the model
alternates between reasoning steps and tool invocations, and every tool
result is appended to a transcript re-sent on the next invocation. Two
properties follow.

\emph{Context accumulates.} Turn $i$'s prompt contains everything from
turns $0\ldots i{-}1$. A program that runs long enough has a prompt of
tens of thousands of tokens, almost all of it prefix already computed on
a previous turn --- which is exactly why serving systems work so hard to
keep it cached~\cite{vllm}.

\emph{Programs are long.} Published measurements of real coding-agent
trajectories report an average of 51 assistant turns per session, with a
tail beyond 200~\cite{polar}. The turn counts we sweep, up to 64, are
therefore at or below what production agents actually do.

\subsection{Autoscaling signals}

Three signals are used in practice. \emph{GPU utilisation} is what
Kubernetes HPA reads by default in every GPU-autoscaling guide.
\emph{Queue depth} --- \code{vllm:num\_requests\_waiting} --- is what the
vLLM production documentation and llm-d guides recommend for KEDA.
\emph{KV-cache utilisation} is the newest, used by llm-d's Workload
Variant Autoscaler and named by HeteroScale as its planned direction.

The first two are compute-side; a parked program registers on neither.
The third is memory-side and does see parked programs --- but as
undifferentiated pressure, without knowing they will return. As we show,
it performs worst of the three.

\subsection{Two Azure traces, two different workloads}

Microsoft released two LLM inference traces from Azure production
services, both November 2023: a \emph{conversation} service and a
\emph{coding} service~\cite{splitwise}. Each record has an arrival
timestamp, a context length, and a generated length. The coding trace is
the closest public proxy for agentic serving: an assistant repeatedly
sends an accumulated buffer and receives a short completion.

\begin{figure}[t]
\centering
\includegraphics[width=\columnwidth]{figs/f1_trace_contrast.pdf}
\caption{The two Azure traces differ in kind, not degree. Measured over
  8{,}819 coding and 19{,}366 chat requests spanning about one hour each.}
\label{fig:traces}
\end{figure}

Figure~\ref{fig:traces} reports three measured differences.

\emph{Shape of the work.} Coding has a mean context of 2{,}048 tokens and
generates a mean of 27.9 --- a ratio of 73:1. Chat is 1{,}155 and 211, a
ratio of 5.5:1. Coding requests are almost pure prefill: the signature of
a loop that accumulates state and emits a short action.

\emph{Burstiness.} Binning arrivals into ten-second windows, coding has an
index of dispersion of 71.3 against 4.0 for chat. Peak-to-mean is 12.8
against 1.8; the inter-arrival coefficient of variation is 13.2 against
1.09. Chat arrivals are nearly Poisson. Coding arrivals are emphatically
not.

The second observation is the one that matters. An autoscaler exists to
track load. A workload with peak-to-mean 1.8 hardly needs one; a workload
with peak-to-mean 12.8 needs one badly. The workload that most requires
autoscaling is precisely the one on which every deployed autoscaler fails.

% =====================================================================
\section{The Parking Problem}
\label{sec:problem}

\subsection{Statement}

Consider a cluster serving agent programs of $T$ turns with mean tool
time $\tau$, and let $s$ be the mean service time of one model
invocation. A program's lifetime is $T(s+\tau)$, of which $T\tau$ is
parked. In steady state with program arrival rate $\lambda_p$, Little's
law gives $\lambda_p T(s+\tau)$ resident programs, of which a fraction
\begin{equation}
\phi \;=\; \frac{\tau}{s+\tau}
\label{eq:parkfrac}
\end{equation}
are parked at any instant. For a coding agent with $s\approx1$\,s and
tool calls of a few seconds, $\phi$ exceeds one half.

Each parked program holds $c_i$ tokens of KV, where $c_i$ is its
accumulated transcript. Total resident KV is $\sum_i c_i$ over running
\emph{and} parked programs, while compute demand counts only the running
ones. The two quantities an autoscaler must reconcile diverge as
$\phi\to1$, and no compute-side metric can see the divergence.

\subsection{Why scale-in is not the inverse of scale-out}

For stateless serving, scaling in is the inverse of scaling out: remove a
replica, drain or retry its in-flight requests, and the cost is bounded
by one request's service time. Under agentic serving it is not. Removing
a replica destroys the KV of every program parked on it, and each pays a
full re-prefill of $c_i$ tokens on its next turn. The cost is
\begin{equation}
\mathcal{C}_{\text{in}} = \sum_{i \in \text{parked}(r)}
\frac{c_i}{R_{\text{prefill}}}
\;+\; \mathbb{1}[\text{reversed}]\cdot t_{\text{cold}}
\label{eq:scaleincost}
\end{equation}
where the second term appears if the decision must later be undone. With
$c_i$ in the tens of thousands and $t_{\text{cold}}$ at 180\,s for a
70B-class model, one bad scale-in costs more than minutes of the replica
it saved.

This asymmetry is why the problem is not solved by tuning a threshold. A
symmetric controller --- and all deployed ones are symmetric --- treats
adding and removing a replica as equal and opposite. Under agentic
serving they are neither.

% =====================================================================
\section{Method}
\label{sec:method}

\subsection{What is real and what is constructed}
\label{sec:construction}

The Azure traces predate agentic serving. They carry no program
identifiers and no tool-call durations, and to our knowledge no public
trace with production arrival times carries both. We state the boundary
explicitly.

\begin{description}
\item[Real, from the trace:] arrival timestamps, and therefore the entire
  arrival process including its burstiness; per-invocation context length;
  per-invocation generated length.
\item[Constructed, and swept:] which invocations compose a program ($T$
  turns), and the tool think-time between turns ($\tau$).
\end{description}

We never fix $T$ and $\tau$ and report one number. We sweep
$T\in\{1,4,16,64\}$ and $\tau\in\{0,8,30\}$\,s and report the surface.
$T{=}1,\tau{=}0$ recovers the original request-level trace exactly; it is
chat, it is our control, and on it the baselines should work --- and do.

The sweep range is calibrated to published measurements of real agent runs
rather than chosen for convenience. Polar reports a mean of 51 assistant
turns per SWE-agent session with a tail past 200~\cite{polar}; our
$T{=}64$ therefore sits near the middle of the real distribution, and
conclusions drawn at $T{=}16$ are conservative.

Context accumulates as in a real agent: turn $i$'s prefill is the whole
transcript so far, growing by what the model generated plus what the tool
returned. Under this construction a 16-turn program reaches a mean final
context of 8{,}510 tokens and a p99 of 34{,}340 --- the right order of
magnitude for deployed coding agents.

\subsection{Simulator}

We simulate a cluster of replicas, each with a KV capacity and a
continuous-batching service model. Aggregate throughput saturates with
batch size: with $k$ concurrent turns each advances at rate
$R/(k+k_{1/2})$, so per-turn latency degrades far more slowly than
$1/k$. This is what batching buys, and a naive processor-sharing model
that missed it would overstate the cost of load and flatter every policy
that over-provisions.

Defaults model a 70B-class model on two A100-80GBs, the regime agentic
deployments run in because agents need strong tool-use models: 62{,}000
tokens of KV per replica (160\,GB minus $\sim$140\,GB of weights at
320\,KB per token), $R = 4{,}000$ token-equivalents/s, $k_{1/2}=16$,
batch $\le 32$, and a 180\,s cold start to load 140\,GB of weights and
warm up. Cold start is swept in \S\ref{sec:eval-cold}.

A program loses its KV two ways: \emph{eviction}, when its replica exceeds
capacity, and \emph{scale-in}, when the autoscaler removes the replica it
is parked on. Either way the next turn re-prefills the whole transcript.
We count the two separately, because attributing the damage is the point.

\subsection{Metrics, and a survivorship trap}
\label{sec:metrics}

Turn service times span two orders of magnitude, so an absolute latency
bound is meaningless. We use \emph{slowdown}: a turn's latency divided by
what it would have taken alone on a replica. The SLO is slowdown $\le3$.

An early version of this study scored SLO attainment over turns that
\emph{completed}, and produced a striking result: our Q-learning
autoscaler achieved perfect attainment. It had learned to shrink the
cluster to one replica, whereupon 861 of 936 programs never finished, and
the handful of short turns that did complete all met the SLO. The metric
rewarded starvation.

We therefore count every turn that was admitted but never completed as a
violation, and report \emph{goodput} --- the fraction of offered programs
that finished --- as a first-class metric alongside SLO. Every number in
this paper uses the corrected definitions. We report the trap because it
is easy to fall into and, where programs are long, it inverts the ranking
of policies.

\subsection{The fairness rule}

Comparing autoscalers on cost alone is trivially gamed by scaling to one,
and on SLO alone by scaling to the maximum. For every workload point we
search for the cheapest \emph{static} cluster meeting the SLO and use it
as the reference: an autoscaler is interesting only if it holds the SLO
for less than a static cluster would cost, because a static cluster is
what you get for not autoscaling at all. All policies are warm-started at
that size and the first 300\,s are excluded, so we measure tracking rather
than an arbitrary initial transient.

\subsection{Policies}

We implement six, spanning all three branches of the standard autoscaling
taxonomy~\cite{gari}. The baselines are what is deployed, not straw men.

\paragraph{Reactive.}
\textbf{HPA (GPU util)} is the Kubernetes autoscaler on
\code{DCGM\_FI\_DEV\_GPU\_UTIL}, target 70\%, with HPA's exact
recommendation formula and the default 300\,s scale-down stabilisation
window. \textbf{KEDA (queue)} scales on
\code{vllm:num\_requests\_waiting}, the metric the vLLM production guides
recommend. \textbf{KV-util} targets 80\% KV-cache utilisation, as llm-d's
Workload Variant Autoscaler does.

\paragraph{Proactive.}
\textbf{Predictive} runs Holt double-exponential smoothing on recent
arrivals and provisions for the forecast one cold-start horizon ahead. It
is a strong baseline precisely because it targets the symptom our results
are full of: reacting too late.

\paragraph{Learning-based.}
\textbf{RL (Q-learning)} is tabular Q-learning over a discretised state
--- replica count, queue depth, GPU utilisation, KV utilisation --- with
actions $\{-2,-1,0,+1,+2\}$ replicas and reward
$-(\text{cost}) - \lambda(\text{SLO violations})$. This is the classic
formulation in the cloud-autoscaling literature~\cite{barrett,gari}. It
trains on a disjoint 1{,}200\,s prefix for 25 episodes and is evaluated on
the suffix; we tuned $\lambda$ over $\{12,40,120,400\}$ and report the
best. \textbf{RL + parked} is the same agent with the parked count added
to its state: the ablation separating ``a better algorithm'' from ``a
better signal''.

\paragraph{Ours.}
\sys changes one thing: occupancy counts parked programs.
\begin{equation}
N = \max\!\Big(
\Big\lceil \tfrac{\text{running}+\text{parked}+\text{queued}}{B\,\theta}\Big\rceil,
\Big\lceil \tfrac{\text{KV used}}{K\,\theta}\Big\rceil \Big)
\label{eq:policy}
\end{equation}
with batch $B$, per-replica KV $K$, and target $\theta$. The second term
makes memory a floor rather than an afterthought; the first counts a
program waiting on a tool as what it is, a seat that is taken. Scale-in
passes through a 60\,s stabilisation window, an order of magnitude shorter
than HPA's default, because the parked term already supplies the damping.

% =====================================================================
\section{Evaluation}
\label{sec:eval}

\subsection{How much of the cluster is parked?}
\label{sec:eval-parked}

\begin{figure}[t]
\centering
\includegraphics[width=0.62\columnwidth]{figs/f2_parked_fraction.pdf}
\caption{Fraction of resident programs parked in a tool call, measured on
  a fixed cluster so no autoscaler can confound it.}
\label{fig:parked}
\end{figure}

Figure~\ref{fig:parked} answers the question motivating everything else.
At $T{=}1$ --- chat --- nothing is parked, by construction. At $T{=}4$,
25.1\% of resident programs are inside a tool call at any instant; at
$T{=}8$, 48.6\%; at $T{=}16$, 43.1\%. Roughly half the cluster's occupancy
is invisible to any compute-side metric, and Eq.~\eqref{eq:parkfrac}
predicts the trend.

\subsection{Every autoscaler degrades}
\label{sec:eval-main}

\begin{figure*}[t]
\centering
\includegraphics[width=\textwidth]{figs/f3_main_slo.pdf}
\caption{SLO attainment (top) and goodput (bottom) across the sweep. At
  $T{=}1$ the workload is chat and the deployed policies work. As programs
  lengthen and tool time grows every baseline collapses; \sys degrades
  most slowly. The Q-learner (grey) completes nothing anywhere.}
\label{fig:main}
\end{figure*}

\begin{table}[t]
\centering\small
\caption{Agentic region ($T\ge4$, $\tau\ge8$\,s), averaged. Cost is
  GPU-seconds relative to the cheapest static cluster meeting the SLO.}
\label{tab:main}
\begin{tabular}{lccc}
\toprule
Policy & SLO attain. & Goodput & Cost \\
\midrule
HPA (GPU util)      & 0.374 & 0.308 & 0.31 \\
KEDA (queue)        & 0.330 & 0.286 & 0.24 \\
KV-util (llm-d)     & 0.253 & 0.165 & 0.23 \\
Predictive (Holt)   & 0.449 & 0.413 & 0.28 \\
RL (Q-learning)     & 0.115 & 0.000 & 0.05 \\
\midrule
\sys (ours)         & \textbf{0.655} & \textbf{0.527} & 0.40 \\
\bottomrule
\end{tabular}
\end{table}

Figure~\ref{fig:main} and Table~\ref{tab:main} are the main result.

On the control point --- $T{=}1$, the unmodified trace, i.e.\ chat --- the
deployed policies are fine: KEDA attains 0.986, HPA 0.897. This matters,
because it establishes that the baselines are correctly implemented and
well suited to what they were designed for. Nothing in our simulator
handicaps them.

Averaged over the agentic region, attainment collapses to 0.253--0.449
and goodput to 0.165--0.413. \sys holds 0.655 and 0.527: $1.46\times$ the
best baseline on SLO and $1.28\times$ on goodput.

The cost column identifies the failure mode. Every baseline spends
0.05--0.31 of the cheapest static cluster's GPU-seconds. They are not
mis-tracking symmetrically; they are systematically \emph{shrinking} the
cluster, because that is what their signals tell them to do, and paying
in SLO and goodput. \sys spends 0.40 --- still less than static, but it is
the only policy that does not read a parked agent as an empty seat.

KV-util deserves comment, being the newest and, on paper, the one that
should see parked programs. It performs worst of the three reactive
policies (0.253 SLO, 0.165 goodput). Seeing memory is not enough: KV
utilisation conflates memory about to be reused with memory merely
occupied, and a replica full of parked programs reads identically to one
that is genuinely saturated.

\subsection{Where the damage comes from}
\label{sec:eval-damage}

\begin{figure}[t]
\centering
\includegraphics[width=\columnwidth]{figs/f4_scalein_damage.pdf}
\caption{Left: programs whose KV cache was destroyed by a scale-in
  decision. Right: the resulting SLO attainment.}
\label{fig:damage}
\end{figure}

Figure~\ref{fig:damage} separates the two ways a program loses its cache.
Averaged over $T\ge4$, scale-in destroys the KV of 3{,}108 programs per
run under KEDA and 422 under HPA, against 221 under \sys --- an order of
magnitude fewer. KEDA is worst precisely because its signal is the most
badly mismatched: an empty queue during a wave of tool calls reads as an
idle cluster, and it removes replicas full of parked state.

Each destroyed program pays a full transcript re-prefill on its next turn.
That recompute lands on a cluster just made smaller, raising latency,
which the autoscaler eventually sees --- 180 seconds later, after a cold
start. Equation~\eqref{eq:scaleincost} prices this loop; the figure
measures it.

\subsection{Why reinforcement learning fails here}
\label{sec:eval-rl}

The Q-learner is the worst policy we evaluated. It attains 0.115 SLO and
completes \emph{zero} programs across the agentic region. Adding the
parked count to its state does not help (0.107 versus 0.115).

The failure is instructive and, we argue, structural rather than a tuning
artefact. In the reward $-(\text{cost}) - \lambda(\text{violations})$, the
cost term is immediate, dense, and exactly attributable to the action just
taken. The violation term is delayed by up to a full cold start: remove a
replica now, and the SLO damage appears 180 seconds and a dozen decisions
later, by which time credit has diffused across the trajectory. The agent
reliably learns the half of the reward it can see.

We tuned $\lambda$ across a factor of 33 and the pathology persisted; at
$\lambda{=}12$ it emerged in the most dramatic form, the one that
motivated the metric correction in \S\ref{sec:metrics}. We do not claim
that no RL autoscaler can work --- a method with a longer effective
horizon, a shaped reward, or a learned model of cold start might well. We
claim the obvious formulation, the one in the standard autoscaling
literature, transfers badly to a setting where the dominant cost is
delayed by minutes, and that this should be reported.

That the parked signal does not rescue it is itself informative. The two
failures are independent: the baselines fail because their signal is
blind, and the Q-learner fails because its credit assignment is broken.
Fixing one does not fix the other.

\subsection{The advantage is a function of cold start}
\label{sec:eval-cold}

\begin{figure}[t]
\centering
\includegraphics[width=\columnwidth]{figs/f5_coldstart.pdf}
\caption{Sweeping replica cold-start time at $T{=}8,\tau{=}8$\,s.}
\label{fig:cold}
\end{figure}

\begin{table}[t]
\centering\small
\caption{SLO attainment against cold-start time ($T{=}8$, $\tau{=}8$\,s).}
\label{tab:cold}
\begin{tabular}{lcccc}
\toprule
Cold start & 30\,s & 90\,s & 180\,s & 300\,s \\
\midrule
HPA         & 0.919 & 0.919 & 0.348 & 0.268 \\
KEDA        & 0.891 & 0.700 & 0.373 & 0.327 \\
Predictive  & 0.824 & 0.713 & 0.672 & 0.581 \\
\sys        & \textbf{0.955} & \textbf{0.836} & \textbf{0.864} & \textbf{0.712} \\
\midrule
\sys / HPA  & 1.04$\times$ & 0.91$\times$ & 2.48$\times$ & 2.66$\times$ \\
\bottomrule
\end{tabular}
\end{table}

Figure~\ref{fig:cold} and Table~\ref{tab:cold} isolate the mechanism. At a
30\,s cold start all policies are within a few percent: destroying state
costs little when you can rebuild it quickly, so carelessness is nearly
free. As cold start grows the baselines fall away --- HPA from 0.919 to
0.268, KEDA from 0.891 to 0.327 --- while \sys degrades gracefully from
0.955 to 0.712.

This is the cleanest evidence that the mechanism we claim is the one
operating. \sys is not a better forecaster; it makes no forecast at all.
It simply declines to destroy state that is expensive to rebuild. Make
rebuilding cheap and its advantage disappears, exactly as that account
predicts.

The predictive baseline is the most interesting comparison. It is the best
baseline at large cold start (0.672 at 180\,s) because forecasting ahead
is the textbook answer to slow provisioning --- and it is still well
behind \sys, because forecasting the \emph{arrival} process says nothing
about the programs already resident and parked. The two are
complementary; combining them is future work we did not do.

\subsection{The effect is not specific to the coding trace}

\begin{figure}[t]
\centering
\includegraphics[width=0.85\columnwidth]{figs/f6_both_traces.pdf}
\caption{Both Azure traces, replayed as chat ($T{=}1$) and as agent
  programs ($T{=}8,\tau{=}8$\,s).}
\label{fig:both}
\end{figure}

Figure~\ref{fig:both} repeats the experiment on the conversation trace. At
$T{=}1$ all policies do well on both. At $T{=}8$ on the chat trace, HPA
falls to 0.689 and KEDA to 0.506 while \sys holds 0.846; on the coding
trace the same comparison is 0.348 and 0.373 against 0.864. The effect is
a property of program structure, not of one trace, and it is worse on the
burstier workload --- which is the one that actually looks like an agent.

% =====================================================================
\section{Threats to Validity}
\label{sec:threats}

\textbf{This is simulation.} No number here was measured on a running GPU
cluster. The service model, cold-start cost, and KV capacity are
parameters chosen to match published hardware characteristics, not
measurements we took. We claim a mechanism and a direction, not a
coefficient. The natural next step, which we have not done, is to replay
these workloads against vLLM on real hardware; we designed the
simulator's metric interface to mirror vLLM's exported Prometheus metrics
so that substitution is mechanical.

\textbf{Program structure is constructed.} The composition of invocations
into programs and the tool think-times are ours, not Azure's, because no
public trace with production arrival times carries them. We mitigate this
three ways: sweeping both parameters over an order of magnitude and
reporting the whole surface rather than a point; calibrating the sweep
range to published statistics of real agent trajectories~\cite{polar}
rather than to convenience; and including $T{=}1$ as a control where the
construction vanishes entirely and the baselines are seen to work. If real
agent programs have a structure our sweep does not contain, our
conclusions may not transfer.

\textbf{The baselines are reimplementations.} We implement HPA's formula
and stabilisation window, KEDA's queue trigger, and llm-d's KV-utilisation
target from their documentation. Production deployments carry tuning we do
not model, and a well-tuned KEDA might do better than ours.

\textbf{Our RL agent may be too weak.} Tabular Q-learning is the classic
formulation, not the strongest possible. We report its failure and our
diagnosis, and explicitly do not claim RL cannot work here.

\textbf{\sys buys SLO with money.} It spends 0.40 of the static reference
while baselines spend 0.05--0.31. Whether that trade is right depends on a
price the operator sets. Our evaluation assumes SLO and goodput are worth
paying for; an operator with the opposite preference should read
Table~\ref{tab:main} in the other direction.

% =====================================================================
\section{Related Work}

\textbf{Agent serving systems} treat the agent program, rather than the
request, as the scheduling unit. Autellix~\cite{autellix} schedules on
cumulative program service time and reports large throughput gains over
vanilla vLLM. Continuum pins KV across tool-call pauses using a predicted
time-to-live. MORI places programs across HBM and DRAM by an idleness
score. CONCUR bounds agent-level concurrency to prevent eviction
thrashing. SAGA schedules whole workflows atomically. Every one optimises
\emph{within a cluster of fixed size}; none asks how large the cluster
should be. That question is ours, and it is not answered by better
scheduling: a perfectly scheduled cluster of the wrong size still misses
its SLO.

\textbf{LLM autoscaling} is the complementary literature.
TokenScale~\cite{tokenscale} scales disaggregated prefill and decode
pools on token velocity. HeteroScale coordinates P/D ratios across
heterogeneous hardware and names KV-aware autoscaling as future work.
ServerlessLLM~\cite{serverlessllm} attacks model-loading cold start
directly, which \S\ref{sec:eval-cold} shows is the single most leveraged
parameter. llm-d's Workload Variant Autoscaler scales on KV utilisation.
All target request-level chat serving. We show that the workload
assumption, not the scaling mechanism, is what breaks.

\textbf{Classical elasticity.} Herbst et al.~\cite{herbst} defined cloud
elasticity as the timeliness and accuracy of resource adaptation, and the
taxonomy of reactive, proactive, and learning-based policies is due to the
surveys~\cite{gari}. The parked-program phenomenon is a case their
framework does not cover: a unit of work that holds resources while
contributing nothing to any utilisation signal. In their vocabulary,
supply and demand are both mismeasured, in opposite directions, by the
same metric.

\textbf{RL for autoscaling} has a long history~\cite{barrett,gari}. Our
negative result is narrow --- one formulation, one workload --- but
identifies a structural reason, that cost is immediate while SLO damage is
delayed by a cold start, which we expect to generalise to any setting
where provisioning latency greatly exceeds the control interval.

% =====================================================================
\section{Conclusion}

Autoscaling LLM clusters is a solved problem for chat and an open one for
agents, and the reason is a single accounting error. A program waiting on
a tool is not idle capacity; it is occupancy that no compute-side metric
can see. Every deployed autoscaler makes this error. On agentic replays of
production Azure traces it costs them roughly half their SLO attainment
and two thirds of their goodput, and it costs a textbook RL autoscaler
everything.

Counting parked programs recovers most of the loss, and does so by
declining to destroy state rather than by predicting anything --- which is
why its advantage tracks the cost of rebuilding that state so precisely.
We release the simulator, the workload builder, and every script that
produced a number in this paper.

\appendix
\section{Sensitivity}
\label{app:sens}

\begin{figure}[h]
\centering
\includegraphics[width=\columnwidth]{figs/f7_sensitivity.pdf}
\caption{Sensitivity at $T{=}8,\tau{=}8$\,s. Left: KV per replica, i.e.\
  model size. Middle: maximum batch. Right: \sys's target $\theta$.}
\label{fig:sens}
\end{figure}

Figure~\ref{fig:sens} sweeps the three parameters most likely to be
questioned.

\textbf{Model size.} Varying KV per replica from 31k tokens (a 140B model
on two A100s) to 500k (an 8B model), \sys stays within a narrow band,
0.752--0.789, while the baselines swing violently: HPA between 0.214 and
0.855, KEDA between 0.439 and 0.909. The baselines are not merely worse on
average; they are unpredictable, because whether their signal happens to
correlate with true occupancy depends on where the binding constraint sits
for that model size. \sys does not depend on that coincidence.

\textbf{Batch size.} \sys is strongest at small batches (0.918 at
$B{=}8$) and weakest at large ones (0.573 at $B{=}64$). This is expected
from Eq.~\eqref{eq:policy}: with a very large batch the compute term
rarely binds and the policy reduces to its memory floor, which is
KV-util with extra steps. Practical vLLM deployments of 70B-class models
run at $B\le64$.

\textbf{Target threshold.} Across $\theta\in[0.5,0.9]$ attainment stays
between 0.778 and 0.869 with no monotone trend, so \sys does not require
careful tuning --- which matters, because an autoscaler that needs
per-workload tuning is one that will be deployed with its defaults.

\bibliographystyle{ACM-Reference-Format}
\begin{thebibliography}{99}
\small
\bibitem{splitwise} P.~Patel et al. Splitwise: Efficient generative LLM
inference using phase splitting. In \emph{ISCA}, 2024. Traces:
\url{https://github.com/Azure/AzurePublicDataset}.
\bibitem{react} S.~Yao et al. ReAct: Synergizing reasoning and acting in
language models. In \emph{ICLR}, 2023.
\bibitem{vllm} W.~Kwon et al. Efficient memory management for large
language model serving with PagedAttention. In \emph{SOSP}, 2023.
\bibitem{autellix} M.~Luo et al. Autellix: An efficient serving engine for
LLM agents as general programs. arXiv:2502.13965, 2025.
\bibitem{tokenscale} R.~Lai et al. TokenScale: Timely and accurate
autoscaling for disaggregated LLM serving with token velocity.
arXiv:2512.03416, 2025.
\bibitem{serverlessllm} Y.~Fu et al. ServerlessLLM: Low-latency serverless
inference for large language models. In \emph{OSDI}, 2024.
\bibitem{herbst} N.~Herbst, S.~Kounev, R.~Reussner. Elasticity in cloud
computing: What it is, and what it is not. In \emph{ICAC}, 2013.
\bibitem{barrett} E.~Barrett, E.~Howley, J.~Duggan. Applying
reinforcement learning towards automating resource allocation and
application scalability in the cloud. \emph{Concurrency and Computation},
25(12), 2013.
\bibitem{gari} Y.~Gar\'i et al. Reinforcement learning-based application
autoscaling in the cloud: A survey. \emph{Engineering Applications of
Artificial Intelligence}, 2021.
\bibitem{polar} Polar: Agentic RL on any harness at scale.
arXiv:2605.24220, 2026. Reports 51 assistant turns per session on average
with a tail beyond 200.
\bibitem{dynamollm} J.~Stojkovic et al. DynamoLLM: Designing LLM
inference clusters for performance and energy efficiency. In \emph{HPCA},
2025.
\bibitem{distserve} Y.~Zhong et al. DistServe: Disaggregating prefill and
decoding for goodput-optimized LLM serving. In \emph{OSDI}, 2024.
\end{thebibliography}

\end{document}
